\begin{abstract}
Radial supporting cuts require a boundary search before linearization, whereas point tangents linearize directly at a violated point. We study whether this extra work pays off in outer approximation of convex generalized disjunctive programs. Both policies use established tangent inequalities, lifted as perspective cuts in hull masters or deactivated in big-$M$ masters, within a common implementation assisted by nonlinear programming (NLP) and matched single- or multiple-tree execution. We state the conditions for valid cuts and finite fixed-residual separation, analyze small selection weights and equivalent row representations, and distinguish these guarantees from numerical objective-gap acceptance. A reproducible study retains 1464 benchmark outcomes on 51 generated controls and 27 external models, together with 420 continuous-cone reference calls. On 33 held-out controls, single-tree extended supporting hyperplanes (ESH) solve all instances with either formulation, versus 32 for extended cutting planes (ECP) with hull and 31 with big-$M$. Across three schedules with unchanged search seeds, common-solved shifted mean wall times favor ESH by 4--6\%. Mean cut counts decrease by 6--18\% across the four primary matched configurations, while mean recorded row-separation time per run increases by factors of 5.5--6.6. External accepted counts are equal, and supported exact conic comparisons favor the conic formulation in common-solved mean time. Removing integer-point NLP recovery reduces pilot successes from 69/72 to 1/72. The evidence supports a modest separator-policy benefit in this implementation, with larger effects from formulation and tree organization; it does not establish a new cut family or general solver superiority.
\end{abstract}
\noindent\textbf{Keywords:} generalized disjunctive programming; convex mixed-integer nonlinear programming; perspective cuts; extended supporting hyperplanes; outer approximation; computational reproducibility.
